\section*{Leitidee}
Bisher haben wir Funktionen vor allem als einzelne Rechenvorschriften betrachtet.  Nun gehen wir einen Schritt weiter: Wir fassen Funktionen selbst als Elemente eines Raumes auf. Funktionen lassen sich addieren und mit Zahlen multiplizieren und bilden damit � ganz analog zu Vektoren � eine lineare Struktur.

Dadurch wird es m�glich, Methoden der linearen Algebra auf Funktionen anzuwenden. Insbesondere k�nnen wir mit Hilfe eines geeigneten Skalarprodukts Begriffe wie L�nge, Abstand und Orthogonalit�t auch f�r Funktionen definieren.

In technischen Anwendungen ist es h�ufig notwendig, eine Funktion m�glichst gut durch einfachere Funktionen zu approximieren.  Die Orthogonalprojektion liefert daf�r ein systematisches Verfahren: Eine Funktion wird auf einen geeigneten Unterraum projiziert und dort durch diejenige Funktion ersetzt, die den kleinsten Abstand besitzt.

Dieses Prinzip der Projektion bildet die Grundlage f�r viele sp�tere Methoden, insbesondere f�r die Fourieranalyse, bei der Funktionen als �berlagerung harmonischer Schwingungen dargestellt werden.
